In this section, we discuss the 
connections between our single-model data attribution setting and 
the standard data attribution problem, 
the computational cost of our \method{} compared to the baselines 
we consider, and finally some limitations of our method.

\paragraph{Relationship between single-model and standard data attribution.}
Recall (from Section \ref{sec:single_model_data_attribution}) that single-model data attribution 
differs from the standard task of predictive data
attribution in that there is no randomness. From a conceptual perspective, 
this means that while standard data attribution is about predicting how a 
{\em new} model would behave if retrained from scratch on a counterfactual
dataset, single-model data attribution is about predicting how a {\em specific}
model would behave if the data had been different.

Technically, single-model data attribution is a more difficult problem than
standard data attribution in the following natural sense: a {\em perfect}
single-model data attribution method can exactly predict the behavior 
of a specific model on any counterfactual dataset. 
By averaging these predictions over multiple specific models 
(each corresponding to a different learning algorithm seed),
we obtain a perfect ``standard'' predictive data attribution method.
The other direction, however, does not hold: it may be possible to construct
a predictive data attribution method that perfectly predicts the 
average behavior of a learning algorithm, but not the behavior
of a specific model.


\paragraph{Computational cost of data attribution.}
One consideration that we have not yet addressed is the computational cost
of {\em building} the data attribution method $\hat{f}$, which varies
per-method and can differ greatly asymptotically for different scenarios.
In what follows, we compare the cost of building \method{},
TRAK, and EK-FAC in a setting with $N$ training samples and $n$ test samples.
\begin{itemize}
  \item TRAK \citep{park2023trak} has a hyperparameter $k$ which 
  controls the accuracy of the method (mechanistically, $k$ 
  determines the number of intermediate checkpoints used by the method).
  TRAK with $k$ averaged models requires computing
  $k \cdot (N + n)$ per-sample gradients, and storing small 
  randomly-projected versions of them. 
  \item EK-FAC has a similar computational structure to TRAK, but 
  instead of just randomly projecting the per-sample gradients, 
  it applies an iterative algorithm with the goal of better 
  approximating the Hessian. 
  In practice, this algorithm seems to dominate 
  its runtime---we refer the reader to \citep{george2018fast} for more details.
  \item \method{} requires a forward and backward pass through the model
    training process for each of the $n$ test samples. The cost 
    of the forward and backward pass scale linearly 
    in the number of train samples $N$, 
    so the full complexity is on the order of $N\cdot n$.
\end{itemize}
In summary, \method{} is faster than TRAK and EK-FAC when $n$ is
small---for example, when $n = 1$ (attributing a single test sample)
on the CIFAR-10 dataset, 
\method{} costs essentially 3-5x as much as training a single model
(15 minutes on a single A100 GPU),
while TRAK and EK-FAC are both 20-100x slower.
As $n$ increases, however, \method{}'s cost increases linearly,
while TRAK and EK-FAC's costs stay roughly constant.




\paragraph{Limitations and future work.}
One prominent limitation is the one discussed in the previous paragraph: 
\method{} becomes expensive as the number of test samples $n$ increases.
The reason for this is that each test sample requires its 
own metagradient calculation, and each of these calculations 
costs 3-5\texttimes{} as much as model training in practice
(see \citet{engstrom2025optimizing} for an in-depth discussion).
An interesting avenue for future work is to develop (a) more efficient 
methods for calculating metagradients, and (b) ways to leverage metagradients to
compute data attribution for multiple test samples at once.

On the algorithmic side, \method{} can only attribute for learning algorithms
that are smooth 
(i.e., as described by \citet{engstrom2025optimizing}, 
``metasmooth''). This and
previous work has identified variations on standard learning algorithms that
are smooth---including for language models,
CLIP~\citep{radford2021learning} models, and standard vision models---but
finding such a variation can take work (see Section 2 of
\citet{engstrom2025optimizing} for a discussion). 
As a simple workaround, we have found
that simply pretraining a model state for a few hundred iterates
(and treating the resulting weights as a ``pretrained'' model initialization) 
is enough to induce smoothness.
Depending on the scenario, this workaround may or may not be acceptable
(i.e., using this method will not allow one to account for the impact of
{\em all} data, but only the data seen after the ``initialization'' stage).
